% ============================================================================
\uxsection{Ubicación de los criterios de evaluación}

La siguiente matriz indica dónde se atiende cada criterio de la rúbrica de la actividad.
Los porcentajes corresponden a los pesos publicados y la modalidad distingue las secciones
elaboradas en equipo de las elaboradas de forma individual por cada integrante.

\begin{uxtabla}{Z{0.9}C{1.3cm}C{2.0cm}Z{1.1}}{Correspondencia entre los criterios de la rúbrica y las secciones del documento}
\uxheadrow \thd{Criterio de la rúbrica} & \thd{Valor} & \thd{Modalidad} & \thd{Dónde se atiende} \\ \midrule
Portada con los nombres de los integrantes & 2\,\% & Equipo & Portada \\
Introducción & 3\,\% & Equipo & Introducción \\
Identificación de la audiencia & 5\,\% & Equipo & Sección 1, apartados 1.1 a 1.7, y Figura 1 \\
Descripción general del problema & 20\,\% & Equipo & Sección 2, apartados 2.1 a 2.5, y Figuras 2 y 3 \\
Definición del producto digital & 20\,\% & Equipo & Sección 3, apartados 3.1 a 3.6, Figura 4 y matriz de trazabilidad de la sección 14 \\
Dos mapas de empatía por cada integrante & 25\,\% & Individual & Secciones 5 a 12: un mapa por persona, agrupados por integrante \\
Dos user personas por cada integrante & 25\,\% & Individual & Secciones 5 a 12: ocho personas, atribuidas por integrante \\
\end{uxtabla}

% ============================================================================
\uxsection{Introducción}

En el sector financiero, o bien, áreas financieras de las corporaciones, la información
necesaria para consultar, validar y analizar datos en muchas ocasiones se encuentra
distribuida en bases y herramientas aisladas. Esta dispersión obliga a los usuarios a
cambiar de sistema, solicitar accesos o apoyarse en conocimiento informal para identificar
la fuente correcta. Como resultado, tareas que deberían ser rápidas pueden convertirse en
procesos largos, repetitivos y difíciles de verificar.

El proyecto propone \textbf{Karisma Data}, un portal centralizado de datos: una plataforma web
modular que concentre el acceso a fuentes financieras institucionales y adapte la experiencia
según la tarea y la responsabilidad de cada persona. La solución considera ocho perfiles de uso que
pueden coexistir o alternarse: consulta operativa, análisis de datos, gobierno y calidad,
control y auditoría, habilitación técnica, supervisión directiva, administración de la
plataforma e integración de aplicaciones.

Este documento presenta una definición inicial de la audiencia, el problema y el producto
digital. También incluye los instrumentos de investigación diseñados para validar las
hipótesis del equipo. La sección individual desarrolla ocho personas y ocho mapas de
empatía correspondientes a esos perfiles.

El documento se organiza en cuatro secciones de equipo y una sección individual. Las
secciones de equipo caracterizan la audiencia (1), el problema (2), el producto digital (3)
y los instrumentos de investigación diseñados para validarlo (4). La sección individual
reúne las ocho personas y sus mapas de empatía, identificadas por integrante. La
caracterización de esta etapa se apoya en investigación documental, en la experiencia
profesional del equipo en áreas que operan con información financiera institucional y en la
observación de canales; su validación con usuarios es el objeto de los instrumentos de la
sección 4. La estructura sigue el caso de estudio guiado del curso (Ramírez Mejía, 2023).

% ============================================================================
\clearpage
\section{Identificación de la audiencia}

La audiencia está formada por personal que consulta, valida, cruza, transforma, publica,
gobierna, audita, integra o supervisa información financiera dentro de una institución. Los
perfiles no son excluyentes: una misma persona puede asumir más de una responsabilidad
según la tarea, aunque cada perfil prioriza objetivos, riesgos y niveles de profundidad
diferentes. La caracterización sigue los componentes del contexto de uso definidos por la
norma ISO 9241-11 (2018): los usuarios, sus tareas y metas, los recursos de que disponen y
los entornos técnico, físico y organizacional en que trabajan.

\subsection{Perfiles principales}

\begin{uxlist}
\lead{Operativo, consulta rápida}{necesita localizar y validar datos puntuales con rapidez.
  Valora una interfaz limpia, un buscador confiable, resultados fáciles de interpretar y
  acceso inmediato a la fuente y fecha de actualización.}
\lead{Analista de datos, profundidad}{necesita combinar fuentes, aplicar filtros complejos,
  revisar metadatos y exportar datos crudos para trabajar en Excel, Python, R, SQL u otras
  herramientas analíticas.}
\lead{Propietario de datos, gobierno y calidad}{es responsable de las definiciones, la
  calidad, la vigencia y las reglas de uso de uno o varios conjuntos de datos. Necesita
  administrar metadatos, atender incidencias, aprobar cambios y comunicar limitaciones de
  forma visible. En ocasiones también requiere acceder a los datos crudos y trabajarlos con
  herramientas analíticas desde un punto de vista gerencial.}
\lead{Riesgos, cumplimiento y auditoría, control y evidencia}{necesita revisar permisos,
  consistencia, trazabilidad y cumplimiento. Busca comprobar quién accedió o modificó
  información, reconstruir el origen de una cifra y documentar excepciones con evidencia
  verificable.}
\lead{Ingeniería y administración de datos, habilitación técnica}{integra fuentes y mantiene
  cargas, esquemas, APIs, permisos técnicos y rendimiento. Necesita monitorear procesos,
  documentar cambios, resolver fallas y ofrecer datos estables a usuarios con distintos
  niveles de experiencia.}
\lead{Directivo, abstracción}{necesita indicadores consolidados, tendencias y señales de
  riesgo para supervisar resultados y tomar decisiones sin recorrer cada base de datos por
  separado.}
\lead{Administración de la plataforma, accesos y continuidad}{gestiona el ciclo de vida de
  las cuentas y los permisos: altas, cambios de rol y bajas. Necesita otorgar el acceso
  mínimo suficiente, revocarlo con la misma rapidez y poder demostrar quién autorizó cada
  permiso y cuándo.}
\lead{Integración de aplicaciones, consumo programático}{construye herramientas internas que
  consumen los datos del portal mediante APIs. Necesita contratos de datos estables y
  documentados, credenciales de autoservicio y aviso anticipado ante cambios de esquema.}
\end{uxlist}

\figuraux{figuras/fig1_mapa_audiencia.png}
  {Mapa de la audiencia: los ocho perfiles según profundidad técnica y responsabilidad sobre
   el dato. Elaboración propia.}{fig:audiencia}

\subsection{Necesidades de la audiencia}

\begin{uxlist}
\lead{Localización}{encontrar la fuente correcta sin depender de contactos informales o
  conocimiento previo del área propietaria.}
\lead{Confianza}{conocer origen, responsable, definición, periodicidad y última
  actualización de cada dato.}
\lead{Velocidad}{resolver consultas simples sin recorrer flujos diseñados para extracciones
  pesadas.}
\lead{Profundidad}{pasar de un resultado general al detalle, filtros y datos crudos cuando
  la tarea lo requiera.}
\lead{Interoperabilidad}{descargar o consumir información en formatos compatibles con sus
  herramientas de trabajo.}
\lead{Continuidad}{guardar búsquedas, filtros, fuentes frecuentes y configuraciones
  personales.}
\lead{Gobierno}{definir propietarios, reglas de uso, niveles de calidad y vigencia de los
  datos.}
\lead{Trazabilidad y control}{reconstruir el origen de una cifra, los accesos, las
  aprobaciones y los cambios realizados.}
\lead{Operabilidad técnica}{monitorear integraciones, cambios de esquema, incidentes y
  rendimiento de los servicios de datos.}
\end{uxlist}

\subsection{Pain points}

\begin{uxlist}
\lead{Información dispersa}{las bases de temas distintos entre ellos, por ejemplo créditos y
  derivados, se encuentran en fuentes separadas.}
\lead{Fricción analítica}{cruzar variables o validar cifras requiere cambiar de sistema,
  formato o responsable.}
\lead{Accesibilidad ineficiente}{los flujos actuales no distinguen entre una consulta rápida
  y una extracción compleja.}
\lead{Silos de conocimiento}{cada área conoce sus propias fuentes, pero tiene poca
  visibilidad de otros datos institucionales útiles.}
\lead{Riesgo de interpretación}{la ausencia de metadatos claros puede provocar que dos
  usuarios comprendan de manera distinta una misma variable.}
\lead{Responsabilidad difusa}{no siempre es claro quién define, corrige o aprueba un dato
  cuando se detecta una inconsistencia.}
\lead{Evidencia limitada}{los accesos, modificaciones y decisiones pueden quedar repartidos
  entre correos, archivos y registros técnicos.}
\lead{Cambios e incidentes poco visibles}{una carga fallida o un cambio de estructura puede
  afectar análisis sin que todos los usuarios se enteren a tiempo.}
\end{uxlist}

\subsection{Características demográficas y profesionales}

Por tratarse de una herramienta institucional, las variables más relevantes no son
únicamente edad o sexo, sino también función, experiencia, nivel técnico, permisos,
frecuencia de uso y responsabilidad sobre los datos. Para el reclutamiento inicial se
propone considerar personal adulto en edad laboral, ubicado principalmente en las sedes
donde operan las áreas financieras o conectado de forma remota. La muestra debe incluir
género, antigüedad, nivel jerárquico y participación como consumidor, propietario,
aprobador, auditor o responsable técnico.

\begin{uxlist}
\lead{Rango de edad orientativo}{21 a 60 años.}
\lead{Ubicación}{Ciudad de México y otras sedes o modalidades remotas donde se consulten
  datos financieros.}
\lead{Experiencia}{desde usuarios con incorporación reciente hasta especialistas con
  conocimiento histórico de las fuentes.}
\lead{Nivel técnico}{básico, intermedio y avanzado; desde consulta en buscador hasta uso de
  SQL, APIs o lenguajes de análisis.}
\end{uxlist}

\subsection{Intereses y fuentes de información}

Los usuarios buscan precisión, trazabilidad, cumplimiento, eficiencia, continuidad y
autonomía. Se relacionan con la información mediante portales internos, bases de datos,
archivos compartidos, correo, mensajería corporativa, documentación técnica, diccionarios,
bitácoras, sistemas de tickets, tableros de monitoreo, reuniones con áreas propietarias y
herramientas de análisis. La investigación deberá identificar cuáles de estos canales
concentran hoy el mayor esfuerzo y cuáles generan más errores, dependencias o riesgos.

\subsection{Método empleado para caracterizar a la audiencia}

La caracterización de esta primera etapa se construyó con tres métodos, sin contacto con
participantes.

\begin{uxlist}
\lead{Investigación documental}{revisión de fuentes públicas sobre el costo de la mala
  experiencia de usuario en herramientas internas y sobre prácticas de gobierno de datos.}
\lead{Experiencia profesional del equipo}{los integrantes trabajan en áreas que consultan,
  validan y administran información financiera institucional, y esa experiencia es la fuente
  de los flujos, el vocabulario y los puntos de fricción descritos.}
\lead{Observación de canales}{revisión de los medios por los que esta audiencia se informa y
  se coordina, que sustenta el apartado 1.5.}
\end{uxlist}

\textbf{\textcolor{uxnavy}{Resultados.}} El método produjo los ocho perfiles de uso descritos
en 1.1, el conjunto de necesidades de 1.2, los ocho pain points de 1.3 y el mapa de canales
de 1.5. El hallazgo central es que los perfiles no son excluyentes: una misma persona
alterna entre consulta rápida y análisis profundo según la tarea, lo que descarta un flujo
único de interfaz y sustenta el principio de revelación progresiva del apartado 3.4.

\begin{uxnota}[Alcance y limitación declarada]
Esta caracterización es una hipótesis de trabajo fundamentada, no un hallazgo empírico. Para
validarla se diseñaron los instrumentos de la sección 4, cuya aplicación inicia en la
siguiente etapa del proyecto y permitirá confirmar, ajustar o descartar cada perfil, así
como completar los indicadores que en 2.2 quedan marcados como pendientes de medición.
\end{uxnota}

\subsection{Contexto de uso}

El diseño centrado en el usuario parte de describir el contexto de uso, es decir, los
usuarios, sus tareas y metas, los recursos de que disponen y los entornos técnico, físico,
social y organizacional en que trabajan (ISO 9241-11, 2018). El mismo conjunto se recuerda
con el acrónimo PACT: personas, actividades, contextos y tecnologías (de Voil, 2020). La
siguiente tabla resume el contexto de uso del portal antes de proponer cualquier solución.

\begin{uxtabla}{L{2.5cm}L{2.1cm}Z{1.15}Z{0.85}}{Contexto de uso del portal según ISO 9241-11 y su consecuencia de diseño}
\uxheadrow \thd{Componente} & \thd{PACT} & \thd{Cómo se manifiesta en el portal} & \thd{Consecuencia de diseño} \\ \midrule
Usuarios & Personas & Ocho perfiles con nivel técnico de básico a avanzado y
  responsabilidades que van de consumir a gobernar el dato; una misma persona cambia de
  perfil según la tarea & Revelación progresiva en lugar de aplicaciones separadas por rol \\
Tareas y metas & Actividades & Localizar y validar una cifra; cruzar fuentes y extraer datos
  crudos; definir y aprobar metadatos; reconstruir accesos y cambios; monitorear cargas;
  supervisar indicadores & Un buscador que resuelve la tarea frecuente y modos avanzados a
  un clic de distancia \\
Recursos & Tecnologías & Navegador y escritorio corporativo, Excel, SQL, Python, herramientas
  de inteligencia de negocio, tickets, correo y mensajería; y recursos expandibles: tiempo,
  esfuerzo y la disponibilidad del colega que conoce la fuente & Exportación e integración
  por API para no romper el flujo de trabajo existente \\
Entornos & Contextos & Técnico: bases heterogéneas con esquemas y accesos separados. Físico:
  oficinas centrales y esquema híbrido, con consulta en escritorio y en tableta. Social y
  organizacional: áreas propietarias, aprobaciones jerárquicas, confidencialidad y auditoría
  & Acceso por permisos y roles, con trazabilidad y bitácora como requisitos, no como extras \\
\end{uxtabla}

Caracterizar el contexto de uso antes de proponer una solución evita diseñar para un usuario
abstracto: cada característica del apartado 3.2 responde a un componente de esta tabla.

% ============================================================================
\clearpage
\section{Descripción general del problema}

\subsection{Definición}

La información financiera institucional está distribuida en fuentes aisladas, con distintos
accesos, estructuras, responsables y niveles de documentación. Esta fragmentación dificulta
localizar, validar, combinar, gobernar, auditar e integrar datos. El problema no es
únicamente tecnológico: también afecta la experiencia de las personas, porque las obliga a
recordar rutas, depender de especialistas, reconstruir evidencia y adaptar su flujo de
trabajo a cada sistema.

El personal que consume, analiza, supervisa, gobierna, controla o habilita información
financiera necesita una forma unificada y confiable de descubrir, comprender, administrar y
utilizar datos, porque la dispersión actual aumenta el tiempo de trabajo, la dependencia
entre áreas y el riesgo de usar información incorrecta, desactualizada o sin trazabilidad
suficiente.

\subsection{Cómo cuantificarlo}

La magnitud del problema se sustenta en tres tipos de evidencia, cada una identificada con
su origen: fuentes públicas verificables sobre el costo de la mala experiencia de usuario en
herramientas internas de empleados, la experiencia profesional del equipo, y los indicadores
que el instrumento de la sección 4 levantará durante el trabajo de campo. Los renglones
marcados como pendientes de medición se enlazan a la pregunta que los resolverá; ninguna
cifra de este documento proviene de una estimación sin fuente.

\begin{uxtabla}{Z{0.9}Z{1.25}C{2.4cm}Z{0.85}}{Cuantificación del problema, con el origen de cada dato}
\uxheadrow \thd{Indicador} & \thd{Qué mide} & \thd{Resultado} & \thd{Origen} \\ \midrule
Costo de la mala experiencia en una herramienta interna & Proyecto de software para personal
  propio descontinuado por dificultad de uso & 125 millones de USD & Maioli (2018), cap. 1 \\
Abandono por diseño o desempeño & Usuarios que dejan de utilizar una aplicación por fallas de
  desempeño & 90\,\% & AppDynamics (2014), como se citó en Maioli (2018) \\
Retorno de invertir en experiencia de usuario & Reducción de errores del usuario y de la
  complejidad percibida & 50\,\% y 65\,\% & On3 (s.\,f.), como se citó en Maioli (2018) \\
Número de fuentes por tarea & Sistemas, carpetas o bases consultadas para completar una
  tarea típica & \textit{Pendiente de medición} & Encuesta 4.1, pregunta \textit{Cantidad de fuentes} \\
Tiempo de localización y validación & Tiempo dedicado a localizar y validar la fuente antes
  de analizar & \textit{Pendiente de medición} & Encuesta 4.1, pregunta \textit{Tiempo} \\
\end{uxtabla}

\subsection{Impacto}

\begin{uxlist}
\lead{Operativo}{aumenta el tiempo de respuesta y dificulta validar datos puntuales durante
  tareas sensibles al tiempo.}
\lead{Analítico}{retrasa cruces, obliga a transformar formatos y reduce el tiempo disponible
  para interpretar resultados.}
\lead{Directivo}{limita la visibilidad consolidada y puede retrasar decisiones o generar
  versiones distintas de una cifra.}
\lead{Gobierno de datos}{dificulta mantener definiciones, responsables, reglas y niveles de
  calidad compartidos.}
\lead{Riesgos y auditoría}{complica reconstruir accesos, cambios y decisiones con evidencia
  suficiente.}
\lead{Tecnológico}{reduce la visibilidad sobre integraciones, incidentes, dependencias y
  cambios de estructura.}
\lead{Institucional}{mantiene silos de conocimiento, duplica esfuerzos y dificulta establecer
  una única fuente de verdad.}
\end{uxlist}

\figuraux{figuras/fig2_impacto_por_perfil.png}
  {Intensidad del impacto del problema por perfil y por consecuencia. Elaboración propia a
   partir del apartado 2.3.}{fig:impacto}

\subsection{Flujo actual y flujo propuesto}

\figuraux{figuras/fig3_flujo_actual_propuesto.png}
  {Flujo de trabajo actual frente al flujo propuesto. Elaboración propia.}{fig:flujo}

\begin{uxtabla}{C{1.1cm}YY}{Comparación paso a paso entre el flujo actual y el propuesto}
\uxheadrow \thd{Paso} & \thd{Flujo actual} & \thd{Flujo propuesto} \\ \midrule
1 & Identificar una posible fuente & Buscar por concepto o necesidad \\
2 & Solicitar acceso o apoyo & Comparar fuentes y metadatos \\
3 & Consultar distintos sistemas & Consultar una vista simple \\
4 & Validar definición y vigencia & Profundizar con filtros cuando sea necesario \\
5 & Transformar o combinar archivos & Exportar o consumir por API \\
6 & Repetir si la fuente no era correcta & Guardar la configuración para reutilizarla \\
\end{uxtabla}

\subsection{El problema en términos de negocio}

Los apartados anteriores describen el problema desde la experiencia de quien trabaja con los
datos. Conviene explicitar por qué además es un problema de negocio. La literatura documenta
que una mala experiencia de usuario en herramientas internas tiene costo directo y medible:
una empresa descontinuó un proyecto de software de 125 millones de dólares tras un piloto,
porque la herramienta resultó demasiado difícil de usar y provocó renuncias en su fuerza de
ventas (Maioli, 2018). Ese costo no aparece en una factura: se distribuye en horas de
trabajo, en decisiones tomadas con información desactualizada y en rotación de personal.

La contraparte también está documentada. Las organizaciones con una experiencia de usuario
efectiva reportan reducciones de hasta 50\,\% en errores del usuario y de 65\,\% en la
complejidad percibida de sus herramientas (Maioli, 2018). Trasladado a este proyecto: cada
minuto que una persona no gasta localizando y validando una fuente es un minuto disponible
para interpretarla, y cada cifra que llega con su origen visible es una discusión que no
ocurre en una reunión directiva.

\begin{uxdestacado}
Por eso el portal se sitúa deliberadamente en la intersección entre las metas institucionales
y las necesidades de las personas, y no en una sola de las dos (Maioli, 2018). Una plataforma
que satisfaga el control institucional pero ignore la fricción cotidiana terminará esquivada
con hojas de cálculo personales, que es exactamente el comportamiento que hoy se busca
corregir.
\end{uxdestacado}

% ============================================================================
\clearpage
\section{Definición del producto digital}

\subsection{Tipo de producto}

\textbf{Karisma Data} es una \textbf{plataforma web interna} de inteligencia y descubrimiento
de datos financieros. Su arquitectura será modular y se adaptará a permisos, responsabilidades
y profundidad de tarea. El portal no sustituye necesariamente todas las bases existentes;
funciona como una capa unificada para descubrirlas, comprenderlas, consultarlas y extraer
información de forma controlada.

El nombre expresa la promesa del producto: un dato que se presenta con su origen, su
responsable y su vigencia a la vista es un dato con autoridad propia, que no necesita ser
defendido en una reunión. La marca acompañará al proyecto en las actividades siguientes.

\subsection{Características principales}

\begin{uxlist}
\lead{Acceso por permisos y roles}{muestra fuentes, funciones y acciones de acuerdo con las
  autorizaciones del usuario.}
\lead{Revelación progresiva}{inicia con resultados simples y permite abrir detalle, filtros o
  herramientas avanzadas según la necesidad.}
\lead{Buscador unificado}{permite buscar por identificador, variable, tema, área, fuente o
  palabra clave.}
\lead{Catálogo de datos}{presenta definiciones, propietarios, periodicidad, calidad,
  actualización y relaciones entre fuentes.}
\lead{Espacios personalizables}{permite guardar módulos, métricas, filtros, fuentes y
  consultas frecuentes.}
\lead{Exportación directa}{ofrece descarga en CSV o Excel y, cuando corresponda, consumo
  mediante APIs.}
\lead{Tableros consolidados}{presenta indicadores y tendencias para seguimiento y toma de
  decisiones.}
\lead{Trazabilidad}{informa origen, fecha de actualización y contexto de uso para reducir
  ambigüedad.}
\lead{Gobierno de metadatos}{permite proponer, revisar y aprobar definiciones, propietarios,
  reglas de uso y niveles de calidad.}
\lead{Flujos de acceso y aprobación}{centraliza solicitudes, autorizaciones y condiciones de
  uso según la sensibilidad de cada fuente.}
\lead{Bitácora y evidencia}{registra accesos, cambios, versiones y decisiones relevantes para
  control y auditoría.}
\lead{Monitoreo técnico}{hace visibles el estado de cargas, APIs, cambios de esquema e
  incidentes que pueden afectar a los consumidores.}
\end{uxlist}

\figuraux{figuras/fig4_concepto_buscador.png}
  {Concepto de la vista de búsqueda: el resultado comunica fuente, propietario, vigencia y
   periodicidad antes de que el usuario tenga que preguntarlas. Elaboración propia.}{fig:concepto}

Las características anteriores se agrupan en siete módulos. La siguiente tabla indica, para
cada uno, el perfil al que sirve de forma primaria y el beneficio concreto que produce, de
modo que ninguna funcionalidad quede sin un usuario y un propósito identificados.

\begin{uxtabla}{L{2.6cm}Z{1.1}L{2.4cm}Z{0.9}}{Módulos del portal, perfil primario y beneficio}
\uxheadrow \thd{Módulo} & \thd{Descripción} & \thd{Perfil primario} & \thd{Beneficio} \\ \midrule
Espacio de trabajo por rol & Pantalla inicial configurable, con los módulos frecuentes y
  valores por defecto según el perfil & Todos & Curva de aprendizaje mínima: la complejidad
  aparece solo cuando la tarea la requiere \\
Buscador y catálogo de datos & Diccionario institucional con búsqueda en lenguaje de negocio,
  fuentes relacionadas y reglas de uso por campo & Operativo y analista & Una sola fuente de
  verdad y fin de la dependencia del colega que conoce la fuente \\
Explorador analítico & Filtros complejos, cruces entre fuentes y vistas tabulares y gráficas
  de alto volumen & Analista de datos & Cruces masivos sin bloquear la pantalla ni transformar
  formatos a mano \\
Tableros y señales de riesgo & Vista consolidada de indicadores con profundización progresiva
  hacia el detalle & Directivo & Comprensión del estado general en un minuto, con el detalle
  disponible pero opcional \\
Asistente conversacional & Consulta en lenguaje natural en la que la herramienta consultada
  queda visible y cada cifra cita su fuente & Todos & Preguntas complejas sin conocer la
  estructura de las bases, con respuesta auditable \\
Exportación e integración & Descargas en CSV y Excel procesadas en segundo plano, con enlace
  de descarga y acceso por API & Analista e integración de aplicaciones & Extracción pesada
  sin fricción y sin bloquear la interfaz \\
Acceso y administración & Autenticación, sesión por rol y panel de altas, cambios de rol y
  desactivación de usuarios & Administración de la plataforma & Gobierno de acceso demostrable
  ante auditoría, con registro de quién autorizó qué \\
\end{uxtabla}

\subsection{Beneficios UX}

\begin{uxlist}
\lead{Eficiencia híbrida}{reduce el tiempo de búsqueda, validación y extracción, y permite
  cambiar de una consulta rápida a un análisis profundo sin perder contexto.}
\lead{Personalización}{disminuye la curva de aprendizaje al mostrar complejidad técnica
  únicamente cuando la tarea la requiere.}
\lead{Estandarización}{facilita el uso de definiciones compartidas y una fuente institucional
  de referencia.}
\lead{Descubrimiento}{hace visibles fuentes relacionadas que el usuario podría no conocer por
  trabajar en otra área.}
\lead{Confianza}{apoya la validación mediante metadatos, responsables y fechas de
  actualización visibles.}
\end{uxlist}

\subsection{Principios iniciales de diseño}

\begin{uxlist}
\lead{La tarea antes que el rol}{la interfaz debe responder a lo que la persona intenta hacer
  en ese momento, porque los modos de uso pueden cambiar.}
\lead{Simplicidad inicial}{la primera vista debe resolver la consulta más frecuente sin
  exponer filtros o términos técnicos innecesarios.}
\lead{Profundidad disponible}{las opciones avanzadas deben estar cerca y ser fáciles de
  descubrir para quien las necesita.}
\lead{Transparencia}{cada resultado debe comunicar origen, vigencia, definición y
  restricciones de uso.}
\lead{Continuidad}{el usuario debe poder guardar, reutilizar y compartir configuraciones de
  manera controlada.}
\end{uxlist}

\subsection{Posicionamiento frente a las alternativas actuales}

La audiencia no parte de cero: hoy resuelve sus tareas con una combinación de bases aisladas,
hojas de cálculo y correo, y en algunos casos con herramientas genéricas de inteligencia de
negocio. La siguiente tabla posiciona la propuesta frente a esas dos alternativas en las
dimensiones que la investigación identificó como críticas.

\begin{uxtabla}{L{2.4cm}Z{1.0}Z{0.95}Z{1.05}}{Posicionamiento frente a las alternativas que la audiencia usa hoy}
\uxheadrow \thd{Dimensión} & \thd{Bases aisladas, hojas de cálculo y correo} &
  \thd{Herramientas genéricas de BI} & \thd{Karisma Data} \\ \midrule
Descubrimiento de datos & Conocimiento no documentado; se pregunta al colega que conoce la
  fuente & Catálogo limitado al espacio de trabajo de la herramienta & Catálogo institucional
  con definición, propietario y reglas de uso, buscable en lenguaje de negocio \\
Modos de uso & Un solo flujo pesado para cualquier necesidad, simple o compleja & Tableros
  fijos, definidos por reporte & Revelación progresiva: consulta rápida, análisis profundo y
  vista directiva en un mismo producto \\
Control de acceso & Carpetas compartidas sin gobierno claro & Licencias y grupos administrados
  por el área técnica & Permisos por perfil, con la matriz de acceso visible y auditable \\
Confianza en el resultado & Se valida comparando contra reportes anteriores o preguntando &
  Respuestas sin origen visible para quien las consume & Cada cifra muestra fuente,
  propietario y fecha de actualización \\
Extracción pesada & Bloquea el equipo o requiere una solicitud al área técnica & Límites de
  filas y exportaciones lentas & Procesamiento en segundo plano con enlace de descarga y
  acceso por API \\
Tiempo hasta el dato & Días cuando la respuesta depende de un tercero & Minutos, siempre que
  el tablero ya exista & Segundos para la consulta frecuente, con la profundidad disponible a
  un clic \\
\end{uxtabla}

El análisis competitivo detallado, con productos concretos del mercado y su evaluación por
perfil, corresponde a una etapa posterior del proyecto. Esta tabla se limita a posicionar la
propuesta frente a las alternativas que la audiencia utiliza hoy, que son las que definen su
punto de comparación real.

\subsection{Cómo se medirá el éxito}

Definir las métricas desde la etapa de concepción evita que la evaluación del producto se
reduzca a una opinión al final del proyecto. El enfoque adoptado sigue el ciclo de definir
métricas, definir usuarios y tareas, medir antes, aplicar los cambios y volver a medir
(Maioli, 2018).

\begin{uxtabla}{L{3.1cm}Z{0.9}Z{1.1}L{2.7cm}}{Métricas de éxito del producto y su forma de medición}
\uxheadrow \thd{Métrica} & \thd{Qué responde} & \thd{Cómo se mide} & \thd{Perfil de referencia} \\ \midrule
Tasa de éxito de la tarea & ¿La persona encuentra la fuente correcta sin ayuda? & Porcentaje
  de tareas completadas sin apoyo externo & Operativo \\
Tiempo hasta el dato validado & ¿Cuánto tarda en confirmar una cifra y su vigencia? &
  Cronometraje de la tarea contra la línea base del flujo actual & Operativo \\
Número de sistemas por tarea & ¿Se redujo la dispersión de fuentes? & Conteo de fuentes
  consultadas antes y después & Analista de datos \\
Solicitudes de apoyo & ¿Disminuyó la dependencia de un colega? & Frecuencia con que la
  persona necesita preguntar para completar la tarea & Operativo y analista \\
Confianza en el resultado & ¿Usaría la cifra sin verificarla en otro lado? & Escala declarada
  al terminar cada tarea & Directivo y propietario de datos \\
Satisfacción de uso & ¿Cómo percibe la usabilidad del portal en conjunto? & Cuestionario
  estandarizado de usabilidad al cierre de la sesión & Todos los perfiles \\
\end{uxtabla}

La línea base de las tres primeras métricas proviene de los indicadores marcados como
pendientes de medición en el apartado 2.2: el mismo instrumento que cuantifica el problema
establece el punto de partida contra el cual se evaluará la solución.

% ============================================================================
\clearpage
\section{Instrumentos de investigación de usuario}

Los instrumentos están orientados a validar el problema, distinguir tipos de tarea y
priorizar funciones. Deben aplicarse a usuarios reales de las áreas involucradas y ajustarse
a las políticas de confidencialidad de la institución. No se solicitarán cifras, clientes ni
datos sensibles durante las sesiones.

El diseño combina deliberadamente dos tipos de indagación. La encuesta y la entrevista son
métodos de opinión: recogen lo que las personas dicen que hacen, información que atraviesa
filtros cognitivos y sociales y que por sí sola es una base frágil para decidir
(de Voil, 2020). La observación contextual es un método empírico: registra lo que las
personas efectivamente hacen en su entorno de trabajo. Los tres instrumentos se complementan,
y por esa razón la entrevista pregunta por hechos recientes y no por intenciones futuras.

\subsection{Encuesta diagnóstica}

\textbf{\textcolor{uxnavy}{Objetivo:}} identificar frecuencia de uso, fuentes consultadas,
esfuerzo, pain points y necesidades de exportación. \textit{Duración estimada: 5 a 7 minutos.}

\begin{uxpreg}
\item \textbf{Perfil laboral.} ¿Cuál es tu área, puesto y antigüedad aproximada en la
  institución? \textit{\textcolor{uxmuted}{Respuesta: abierta breve.}}
\item \textbf{Frecuencia.} ¿Con qué frecuencia consultas información financiera?
  \textit{\textcolor{uxmuted}{Respuesta: varias veces al día, diariamente, semanalmente,
  mensualmente o con menor frecuencia.}}
\item \textbf{Tipos de tarea.} ¿Qué tareas realizas con mayor frecuencia?
  \textit{\textcolor{uxmuted}{Respuesta: consulta puntual, validación, cruce de variables,
  análisis, exportación, supervisión, publicación o mantenimiento, gobierno de metadatos,
  aprobación de accesos, auditoría, integración o monitoreo técnico, u otra.}}
\item \textbf{Cantidad de fuentes.} ¿Cuántas fuentes, sistemas o carpetas consultas
  normalmente para completar una tarea? \textit{\textcolor{uxmuted}{Respuesta: rango numérico.}}
\item \textbf{Tiempo.} En una tarea típica, ¿cuánto tiempo dedicas a localizar y validar la
  fuente antes de analizar la información? \textit{\textcolor{uxmuted}{Respuesta: rango de minutos.}}
\item \textbf{Facilidad.} ¿Qué tan fácil es encontrar la fuente correcta?
  \textit{\textcolor{uxmuted}{Respuesta: escala de 1 a 5.}}
\item \textbf{Dependencia.} ¿Con qué frecuencia necesitas pedir ayuda a otra persona o área
  para localizar, interpretar o acceder a un dato?
  \textit{\textcolor{uxmuted}{Respuesta: escala de frecuencia.}}
\item \textbf{Dificultades.} ¿Cuáles son las principales dificultades que enfrentas?
  Selecciona hasta tres. \textit{\textcolor{uxmuted}{Respuesta: acceso, búsqueda,
  definiciones, actualización, formato, rendimiento, permisos, desconocimiento de fuentes u otra.}}
\item \textbf{Confianza.} ¿Qué información necesitas ver para confiar en un dato?
  \textit{\textcolor{uxmuted}{Respuesta: propietario, definición, actualización, periodicidad,
  calidad, fuente original u otra.}}
\item \textbf{Formatos.} ¿Qué formatos utilizas para trabajar fuera del sistema?
  \textit{\textcolor{uxmuted}{Respuesta: CSV, Excel, API, conexión a base, PDF u otro.}}
\item \textbf{Herramientas.} ¿Qué herramientas utilizas después de extraer los datos?
  \textit{\textcolor{uxmuted}{Respuesta: Excel, Power BI, Python, R, SQL, SAS u otra.}}
\item \textbf{Personalización.} ¿Qué tan útil sería guardar búsquedas, filtros o fuentes
  frecuentes? \textit{\textcolor{uxmuted}{Respuesta: escala de 1 a 5.}}
\item \textbf{Prioridad.} Si existiera un portal único, ¿qué tarea te gustaría resolver
  primero en él? \textit{\textcolor{uxmuted}{Respuesta: abierta.}}
\item \textbf{Experiencia reciente.} Describe una ocasión en la que encontrar o validar un
  dato tomó más tiempo de lo esperado. \textit{\textcolor{uxmuted}{Respuesta: abierta.}}
\end{uxpreg}

\subsection{Guía de entrevista semiestructurada}

\textbf{\textcolor{uxnavy}{Objetivo:}} comprender el flujo real de una tarea, las decisiones
del usuario, sus atajos y los criterios con los que evalúa la confiabilidad.
\textit{Duración estimada: 30 a 40 minutos.} El guion privilegia preguntas de recuerdo sobre
hechos recientes en lugar de preguntas hipotéticas, y evita preguntas dirigidas, siguiendo la
taxonomía de Portigal (2013).

\begin{uxpreg}
\item \textbf{Contexto.} Cuéntame sobre tu función y el tipo de información financiera que
  utilizas.
\item \textbf{Responsabilidad sobre los datos.} En tu trabajo habitual, ¿actúas principalmente
  como consumidor, analista, supervisor, propietario, aprobador, auditor o responsable
  técnico? ¿Cuándo cambia esa responsabilidad?
\item \textbf{Tarea reciente.} Piensa en una tarea reciente. ¿Qué necesitabas lograr y por
  qué era importante?
\item \textbf{Punto de partida.} ¿Cómo decidiste dónde buscar primero?
\item \textbf{Flujo.} ¿Qué pasos seguiste desde que inició la necesidad hasta que obtuviste un
  resultado utilizable?
\item \textbf{Fricción.} ¿En qué momento surgieron dudas, retrasos o solicitudes de apoyo?
\item \textbf{Validación.} ¿Cómo verificaste que la fuente, la definición y la fecha eran
  correctas?
\item \textbf{Continuidad.} ¿Qué herramientas o archivos utilizaste después de encontrar la
  información?
\item \textbf{Esfuerzo.} ¿Qué parte del proceso consideras más repetitiva o frustrante?
\item \textbf{Confianza.} ¿Qué información debería mostrar un resultado para que pudieras
  usarlo con confianza?
\item \textbf{Profundidad.} ¿Qué debería ocurrir en una consulta rápida y qué debería
  reservarse para un modo avanzado?
\item \textbf{Riesgos.} ¿Qué riesgos percibes en un portal centralizado? Considera permisos,
  interpretación, rendimiento o dependencia.
\item \textbf{Cambios e incidentes.} ¿Cómo te enteras de una actualización, una falla de carga
  o un cambio en la estructura de una fuente? ¿Qué haces después?
\item \textbf{Cambio prioritario.} Si pudieras cambiar una sola parte del proceso actual,
  ¿cuál sería y por qué?
\end{uxpreg}

\subsection{Guía de observación contextual}

\textbf{\textcolor{uxnavy}{Objetivo:}} contrastar lo que las personas reportan con lo que
efectivamente hacen al resolver una tarea real con datos. \textit{Duración estimada: de 45 a
60 minutos por sesión}, en el puesto de trabajo del participante y sobre una tarea que ya
tenía prevista realizar. Las notas siguen el esquema AEIOU, que ordena la observación en
cinco dimensiones (de Voil, 2020).

\begin{uxlist}
\lead{Actividades}{qué hace la persona, en qué orden, qué pasos repite y cuáles resuelve con
  un atajo propio en lugar del flujo previsto.}
\lead{Entornos}{cómo es el espacio físico y social en que trabaja: interrupciones, pantallas
  disponibles, cercanía de colegas a quienes puede consultar sin costo.}
\lead{Interacciones}{a quién consulta, por qué canal y en qué momento del flujo aparece la
  dependencia de otra persona o de otra área.}
\lead{Objetos}{qué artefactos intervienen: archivos, plantillas, capturas, correos y hojas de
  control. Interesan sobre todo los que trasladan información de un área a otra.}
\lead{Usuarios}{qué caracteriza a la persona observada: experiencia, nivel técnico y los
  criterios con los que decide que un dato es confiable.}
\end{uxlist}

\begin{uxnota}[Regla de registro]
Se anota por separado lo observado y lo interpretado. Una observación es un hecho verificable
—abrió cuatro pestañas, escribió a un colega, comparó contra un archivo anterior—, mientras
que una interpretación es una inferencia del equipo (de Voil, 2020). Al cierre de cada sesión
las interpretaciones se contrastan con el participante, mientras aún se está en su entorno de
trabajo.
\end{uxnota}

\subsection{Plan inicial de aplicación}

\begin{uxtabla}{L{2.4cm}L{2.8cm}YY}{Plan de aplicación de los tres instrumentos}
\uxheadrow \thd{Método} & \thd{Muestra sugerida} & \thd{Cobertura} & \thd{Evidencia esperada} \\ \midrule
Encuesta & 15 a 30 participantes & Los ocho perfiles de uso & Patrones de frecuencia, esfuerzo
  y prioridad \\
Entrevista & 5 a 8 participantes & Al menos cuatro perfiles, con responsabilidades distintas &
  Flujos, lenguaje, decisiones y puntos de dolor \\
Observación & 3 a 5 tareas & Consulta, análisis, gobierno, control e integración & Pasos
  reales, cambios de sistema y dependencias \\
\end{uxtabla}

\subsection{Criterios de análisis}

Las respuestas se agruparán por tipo de tarea y responsabilidad sobre los datos, no solo por
puesto. El análisis buscará patrones de frecuencia, tiempo, dependencia, confianza,
profundidad, formato de salida, riesgo y necesidad de evidencia. También se registrarán
diferencias entre lo que una persona dice que hace y lo observado durante una tarea real.

\textbf{\textcolor{uxnavy}{Técnica de síntesis.}} Los hallazgos de las tres fuentes se
procesan mediante diagramación de afinidad (de Voil, 2020): cada hallazgo se registra como una
nota independiente y las notas se agrupan por afinidad hasta que los temas emergen del
material, en lugar de partir de categorías definidas de antemano. Los grupos resultantes
alimentan de forma directa los mapas de empatía de este documento y, en la siguiente
actividad, los journey maps por perfil.

% ============================================================================
\clearpage
\aportacion{Jacqueline Sarmiento Cervantes}{Matrícula A01795863 · Personas 1 y 2}

\begin{uxnota}[Nota metodológica]
Las ocho personas son arquetipos hipotéticos en el sentido de Cooper, citado en Hartson y
Pyla (2012): representaciones construidas a partir de la caracterización de la audiencia de la
sección 1, de la experiencia profesional del equipo en el sector financiero y de la literatura
de diseño centrado en el usuario. No corresponden a individuos reales y sus fotografías fueron
generadas con inteligencia artificial. Los mapas de empatía siguen el formato de cuatro
cuadrantes de Osterwalder y Pigneur (2010) con las etiquetas solicitadas. Los instrumentos de
la sección 4 permitirán recalibrar unas y otros con evidencia de campo en la siguiente etapa.
\end{uxnota}

% ---------------------------------------------------------------- Persona 1
\section{Persona 1 · Laura Méndez}

\personahead{imagenes/persona_01_laura.png}{Laura Méndez}
  {“Necesito confirmar el dato y seguir con mi trabajo, no aprender cómo está organizada cada
   base.”}

\begin{personadatos}
Edad & 34 años \\
Sexo & Mujer \\
Ubicación & Ciudad de México \\
Ocupación & Especialista de Operaciones Financieras \\
Experiencia & 8 años en procesos financieros; 3 años en su función actual \\
Nivel técnico & Intermedio: Excel avanzado y consulta de sistemas internos \\
Modo principal & Operativo: consulta rápida y validación puntual \\
\end{personadatos}

\personabloque{Antecedentes}
Laura estudió Administración Financiera y trabaja en un área que atiende solicitudes internas
con tiempos definidos. Conoce bien los procesos de su equipo, pero no siempre sabe qué área es
propietaria de una variable nueva. Tiene una agenda fragmentada entre consultas, seguimiento y
validaciones. Fuera del trabajo le interesan la organización personal, ver series y realizar
actividades con su familia. Vive con su pareja y su hija de seis años en el oriente de la
ciudad, y organiza su jornada para alcanzar la salida de la escuela.

\personabloque{Objetivos}
\begin{uxlist}
\lead{Rapidez}{encontrar una cifra puntual y confirmar su vigencia en pocos minutos.}
\lead{Seguridad}{responder solicitudes con una fuente identificable y una definición clara.}
\lead{Autonomía}{resolver tareas frecuentes sin depender de otra área.}
\lead{Continuidad}{guardar consultas habituales para no reconstruirlas cada vez.}
\end{uxlist}

\personabloque{Pain points y desafíos}
\begin{uxlist}
\lead{Rutas cambiantes}{recuerda nombres de carpetas o sistemas, pero los accesos y
  ubicaciones pueden cambiar.}
\lead{Definiciones ambiguas}{encuentra variables parecidas sin saber cuál corresponde a su
  solicitud.}
\lead{Dependencia}{debe escribir a colegas para confirmar responsable, fecha o fuente y
  también para solicitar actualización de datos.}
\lead{Sobrecarga}{los filtros técnicos y campos avanzados dificultan una consulta que debería
  ser simple.}
\end{uxlist}

\personabloque{Comportamientos y hábitos}
\begin{uxlist}
\lead{Inicio}{busca primero en fuentes conocidas o reutiliza un archivo anterior.}
\lead{Validación}{compara la cifra con reportes previos y pregunta a un contacto de confianza.}
\lead{Herramientas}{trabaja principalmente con navegador, Excel, correo y mensajería
  corporativa.}
\lead{Preferencia}{elige resultados claros, con fecha y fuente visibles, antes que una vista
  con muchas opciones.}
\end{uxlist}

\personabloque{Escenario de uso}
Recibe una solicitud para confirmar una cifra de crédito antes de una reunión. Abre el portal,
escribe el concepto, revisa la fecha de actualización y la definición, filtra por periodo y
copia el dato con su referencia. Si necesita más contexto, abre el detalle sin abandonar la
consulta.

\subsection{Mapa de empatía de Laura}
\mapaempatia
  {\item Solo necesito saber cuál es el dato correcto.
   \item Siempre termino preguntando quién actualiza esta fuente.
   \item Quiero ver la fecha antes de usar el resultado.}
  {\item ¿Esta variable significa lo mismo que en el reporte anterior?
   \item No debería tomar tanto tiempo una validación sencilla.
   \item Si uso la fuente equivocada, tendré que explicar el error.}
  {\item Reutiliza archivos y consultas anteriores.
   \item Compara resultados antes de responder.
   \item Contacta a colegas cuando no identifica al propietario.
   \item Guarda enlaces y notas personales.}
  {\item Presión cuando la solicitud es urgente.
   \item Duda ante nombres o definiciones similares.
   \item Alivio cuando la fuente muestra contexto y vigencia.
   \item Frustración al repetir pasos.}

% ---------------------------------------------------------------- Persona 2
\clearpage
\section{Persona 2 · Diego Hernández}

\personahead{imagenes/persona_02_diego.png}{Diego Hernández}
  {“No me basta con descargar un archivo; necesito entender de dónde viene y poder repetir el
   proceso.”}

\begin{personadatos}
Edad & 29 años \\
Sexo & Hombre \\
Ubicación & Ciudad de México, modalidad híbrida \\
Ocupación & Analista de Datos Financieros \\
Experiencia & 5 años en análisis, automatización y visualización de datos \\
Nivel técnico & Avanzado: SQL, Python, Excel y herramientas de inteligencia de negocio \\
Modo principal & Analista: exploración, cruce, extracción y reutilización \\
\end{personadatos}

\personabloque{Antecedentes}
Diego estudió Actuaría y participa en proyectos que integran información de distintas áreas
financieras. Prefiere procesos reproducibles y documentados. Puede explorar datos con
autonomía, pero pierde tiempo cuando las fuentes utilizan estructuras, nombres o periodos
diferentes. Le interesan la automatización, la visualización, la tecnología y el aprendizaje
continuo. Vive con su pareja en un departamento cercano a la oficina, sin hijos, y dedica
buena parte de su tiempo libre a proyectos personales de programación.

\personabloque{Objetivos}
\begin{uxlist}
\lead{Integración}{descubrir fuentes relacionadas y cruzarlas con criterios consistentes.}
\lead{Reproducibilidad}{guardar consultas, filtros y versiones para repetir el análisis.}
\lead{Acceso}{obtener datos crudos mediante CSV, Excel, conexión o API, según el volumen.}
\lead{Calidad}{identificar cobertura, periodicidad, propietario y limitaciones antes de
  modelar.}
\end{uxlist}

\personabloque{Pain points y desafíos}
\begin{uxlist}
\lead{Estructuras incompatibles}{las fuentes usan identificadores, catálogos o granularidades
  diferentes.}
\lead{Documentación dispersa}{las definiciones se encuentran en correos, archivos o
  conocimiento de especialistas.}
\lead{Exportaciones limitadas}{algunos flujos no funcionan bien con archivos pesados o
  requieren pasos manuales.}
\lead{Trazabilidad}{es difícil reconstruir qué versión, filtro o transformación produjo un
  resultado.}
\end{uxlist}

\personabloque{Comportamientos y hábitos}
\begin{uxlist}
\lead{Exploración}{revisa metadatos y muestras antes de descargar un conjunto completo.}
\lead{Preparación}{documenta transformaciones y conserva scripts o consultas reutilizables.}
\lead{Herramientas}{navegador, SQL, Python, Excel y Power BI.}
\lead{Colaboración}{consulta a propietarios de datos cuando detecta valores atípicos o cambios
  de estructura.}
\end{uxlist}

\personabloque{Escenario de uso}
Debe analizar la relación entre créditos, liquidez y variables de derivados. Busca una fuente
inicial en el catálogo, revisa el diccionario y descubre conjuntos relacionados. Compara
granularidad y fechas, configura filtros, prueba una muestra y después genera una extracción
por API. Guarda la consulta y comparte la referencia para que el análisis pueda repetirse.

\subsection{Mapa de empatía de Diego}
\mapaempatia
  {\item Necesito saber qué representa cada campo.
   \item La descarga debe poder repetirse con los mismos filtros.
   \item Prefiero probar una muestra antes de extraer todo.}
  {\item ¿Las fuentes tienen la misma granularidad y periodo?
   \item ¿Hubo un cambio de estructura que no está documentado?
   \item ¿Cómo explicaría el origen de este resultado a otra persona?}
  {\item Explora metadatos y relaciones entre fuentes.
   \item Construye consultas y scripts reutilizables.
   \item Valida valores atípicos con el área propietaria.
   \item Documenta filtros, fechas y transformaciones.}
  {\item Curiosidad al descubrir fuentes relacionadas.
   \item Frustración con formatos y definiciones inconsistentes.
   \item Desconfianza cuando falta trazabilidad.
   \item Control cuando el proceso es reproducible.}

% ============================================================================
\clearpage
\aportacion{Alexandro Mayoral Terán}{Matrícula A01795899 · Personas 3, 4, 5 y 6}

% ---------------------------------------------------------------- Persona 3
\section{Persona 3 · Roberto Valdez}

\personahead{imagenes/persona_03_roberto.png}{Roberto Valdez}
  {“Si publican un análisis con datos desactualizados o mal interpretados, el problema y las
   explicaciones recaen en mi área.”}

\begin{personadatos}
Edad & 52 años \\
Sexo & Hombre \\
Ubicación & Ciudad de México \\
Ocupación & Subgerente de Información de Estados Financieros y Crédito \\
Experiencia & 18 años en el Banco; experto en catálogos y metodologías crediticias \\
Nivel técnico & Intermedio: bases de datos relacionales, validación de reglas de negocio y
  Excel avanzado \\
Modo principal & Propietario de datos: gobierno, calidad, definiciones y autorización de
  accesos \\
\end{personadatos}

\personabloque{Antecedentes}
Roberto es contador y dirige el equipo que recibe, limpia y clasifica la información
crediticia que envían los bancos. Conoce de memoria los catálogos regulatorios y es muy
meticuloso. Constantemente lidia con solicitudes de otras Direcciones Generales que necesitan
sus bases de datos para investigaciones macroeconómicas, pero que a menudo no comprenden la
complejidad o las reglas detrás de la clasificación de los créditos. Está casado, tiene dos
hijos en edad universitaria y es el principal sostén económico de su casa. Fuera del trabajo
disfruta la lectura de historia y el ajedrez.

\personabloque{Objetivos}
\begin{uxlist}
\lead{Control de calidad}{asegurar que toda el área institucional utilice la definición
  oficial y la versión más reciente de los datos de crédito.}
\lead{Gobierno eficiente}{documentar los metadatos en un solo lugar para evitar responder las
  mismas dudas por correo.}
\lead{Gestión de accesos}{tener la capacidad de aprobar o rechazar quién puede consultar el
  nivel más granular y confidencial de su información.}
\end{uxlist}

\personabloque{Pain points y desafíos}
\begin{uxlist}
\lead{Consultas repetitivas}{pierde horas a la semana explicando a analistas de otras áreas
  qué significa una columna específica en su base de datos.}
\lead{Uso de datos obsoletos}{otras áreas siguen trabajando con extracciones descargadas meses
  atrás y no se enteran cuando la definición o la clasificación oficial ya cambió.}
\lead{Descontrol de versiones}{otras áreas descargan su información, la guardan localmente y
  meses después publican reportes internos con datos viejos y sin trazabilidad.}
\lead{Trazabilidad}{le frustra que no haya un repositorio único donde él pueda decir “esta es
  la fuente de verdad certificada por mi subgerencia”.}
\end{uxlist}

\personabloque{Comportamientos y hábitos}
\begin{uxlist}
\lead{Documentación manual}{mantiene diccionarios de datos detallados, pero actualmente los
  comparte en PDF o Excel por correo institucional.}
\lead{Autorización}{revisa detalladamente el motivo de consulta de cualquier usuario antes de
  solicitar a TI que le den acceso a sus bases.}
\lead{Herramientas}{correo electrónico institucional, sistemas core del banco, Excel y
  software de gestión de requerimientos.}
\end{uxlist}

\personabloque{Escenario de uso}
Roberto necesita actualizar el catálogo de clasificación de cartera de crédito. Sube la nueva
estructura al portal, actualiza las definiciones en el diccionario de datos integrado y marca
la versión anterior como deprecada. El portal notifica automáticamente a todos los analistas
que tienen guardada esa fuente en sus perfiles, asegurando que nadie use metodologías viejas
en su próximo análisis económico.

\subsection{Mapa de empatía de Roberto}
\mapaempatia
  {\item Por favor revisen el diccionario de datos que les mandé por correo antes de cruzar
     variables.
   \item Esa cifra no cuadra porque están usando información desactualizada.
   \item ¿Para qué investigación específica necesitan este nivel de detalle?}
  {\item Ojalá hubiera un lugar donde la gente pudiera leer las reglas del dato sin tener que
     llamarme.
   \item Me preocupa que se filtre información de crédito que aún no es pública.
   \item Tengo que proteger la reputación técnica de mi subgerencia.}
  {\item Dedica tiempo a explicar metodologías por teléfono institucional.
   \item Mantiene un control estricto, casi manual, de quién tiene sus bases de datos.
   \item Revisa los reportes de otras áreas para asegurar que no malinterpretaron sus cifras.}
  {\item Responsabilidad absoluta sobre la veracidad de la información de crédito.
   \item Frustración por ser un cuello de botella para los analistas.
   \item Celoso de la información: valora mucho la confidencialidad y el control.}

% ---------------------------------------------------------------- Persona 4
\clearpage
\section{Persona 4 · Elena Ruiz}

\personahead{imagenes/persona_04_elena.png}{Elena Ruiz}
  {“No puedo basar mi auditoría en cadenas de correos; necesito ver la bitácora de quién tocó
   ese dato y cuándo.”}

\begin{personadatos}
Edad & 36 años \\
Sexo & Mujer \\
Ubicación & Ciudad de México \\
Ocupación & Especialista en Auditoría de Tecnología de Información \\
Experiencia & 7 años auditando sistemas, controles de acceso y procesos financieros críticos \\
Nivel técnico & Avanzado: lectura de bitácoras, SQL y marcos de seguridad \\
Modo principal & Riesgos y auditoría: control, trazabilidad, permisos y recopilación de
  evidencia \\
\end{personadatos}

\personabloque{Antecedentes}
Elena trabaja en la Unidad de Auditoría que reporta a la Junta de Gobierno. Su trabajo no es
analizar la economía, sino vigilar que los sistemas y procesos que manejan los datos del Banco
Central cumplan con las estrictas normativas de seguridad, confidencialidad y transparencia.
Es muy estructurada, analítica y su trabajo requiere que sea escéptica por naturaleza. Es
contadora pública con certificación en auditoría de sistemas. Es madre de un niño de nueve
años y cuida mucho su tiempo, por lo que evita las jornadas extendidas; fuera del trabajo
practica natación y toma cursos de actualización normativa.

\personabloque{Objetivos}
\begin{uxlist}
\lead{Trazabilidad}{poder reconstruir exactamente el ciclo de vida de un dato crítico: quién
  lo subió, quién lo modificó y quién lo descargó.}
\lead{Transparencia de accesos}{verificar de forma rápida qué usuarios tienen privilegios de
  lectura sobre información sensible antes de su publicación.}
\lead{Evidencia inmutable}{obtener reportes del sistema que sirvan como prueba fehaciente para
  sus informes de auditoría, sin depender de testimonios verbales.}
\end{uxlist}

\personabloque{Pain points y desafíos}
\begin{uxlist}
\lead{Fragmentación de bitácoras}{cuando audita un proceso, la evidencia está repartida entre
  sistemas antiguos, correos electrónicos y tickets de mesa de ayuda, haciendo casi imposible
  armar el rompecabezas.}
\lead{Cajas negras}{le incomodan los procesos donde la información se transforma en el Excel
  personal de un analista, perdiendo todo el rastro de la fuente original.}
\lead{Lentitud en la recolección}{depende de que el área de TI le extraiga las bitácoras, lo
  cual retrasa sus tiempos de entrega de la auditoría.}
\end{uxlist}

\personabloque{Comportamientos y hábitos}
\begin{uxlist}
\lead{Recolección de pruebas}{documenta cada paso con capturas de pantalla, oficios y reportes
  extraídos directamente de los sistemas.}
\lead{Análisis de permisos}{cruza constantemente el organigrama con las matrices de acceso a
  los sistemas.}
\lead{Herramientas}{software de auditoría interna, bases de datos, Excel y sistemas de tickets
  institucionales.}
\end{uxlist}

\personabloque{Escenario de uso}
Elena está auditando el proceso de publicación del reporte de Agregados Monetarios. En lugar
de pedir oficios y correos, entra al portal centralizado. Entra al módulo de auditoría, busca
la variable específica y descarga un reporte de trazabilidad que muestra la fecha exacta en
que se cargó el dato, que Roberto Valdez aprobó su calidad, y la lista de los cinco analistas
que tuvieron acceso previo a la publicación. Genera su evidencia en minutos.

\subsection{Mapa de empatía de Elena}
\mapaempatia
  {\item Necesito la evidencia documentada de que el dueño del dato autorizó este acceso.
   \item Esta información estuvo en un archivo compartido, ¿quién más lo vio?
   \item Muéstrenme la bitácora del sistema, no el correo.}
  {\item Si hay una fuga de información y no podemos rastrearla, la Junta de Gobierno nos
     exigirá respuestas.
   \item Es increíble que sigan manejando datos confidenciales fuera del sistema central.
   \item Quiero reducir el tiempo que tardo en juntar pruebas para mi reporte.}
  {\item Levanta tickets a TI pidiendo extracciones de bases de datos.
   \item Entrevista a los usuarios para entender por qué tomaron un atajo en el sistema.
   \item Cruza listas de acceso con la nómina para detectar usuarios inactivos con permisos.}
  {\item Escepticismo natural ante procesos manuales.
   \item Presión por entregar reportes de auditoría sin errores ni suposiciones.
   \item Tranquilidad cuando un sistema habla por sí solo con sus bitácoras.}

% ---------------------------------------------------------------- Persona 5
\clearpage
\section{Persona 5 · Jorge Mendieta}

\personahead{imagenes/persona_05_jorge.png}{Jorge Mendieta}
  {“Mi trabajo es que los fierros no fallen. Si un área cambia la estructura de su base de
   datos y no me avisa, rompen las integraciones de todo el banco.”}

\begin{personadatos}
Edad & 42 años \\
Sexo & Hombre \\
Ubicación & Ciudad de México \\
Ocupación & Administrador de Data Warehouse en la Gerencia de Sistemas de Información \\
Experiencia & 15 años en arquitectura de datos, procesos de extracción y carga, y modelado
  dimensional \\
Nivel técnico & Avanzado: SQL, herramientas de integración y modelado de datos \\
Modo principal & Ingeniería y administración: integración técnica, gobierno técnico y
  rendimiento del almacén de datos \\
\end{personadatos}

\personabloque{Antecedentes}
Jorge trabaja mano a mano con la Gerencia de Información. Mientras que la gerencia define el
significado económico de los datos, Jorge diseña las tuberías técnicas y la estructura del
almacén de datos institucional. Su principal responsabilidad es asegurar que la información
que llega de los bancos se integre limpiamente en las bases históricas, estructurándola de
forma que los analistas puedan consultarla rápidamente. Es ingeniero en sistemas
computacionales con maestría en tecnologías de información. Está casado y tiene dos hijas
pequeñas; procura no llevarse el trabajo a casa, aunque las cargas nocturnas se lo dificultan.
Fuera de la oficina arma equipos de cómputo y sigue el futbol.

\personabloque{Objetivos}
\begin{uxlist}
\lead{Integridad estructural}{asegurar que el almacén de datos sea una única fuente de verdad
  técnica, sin datos duplicados ni inconsistencias.}
\lead{Rendimiento en consultas}{diseñar las tablas para que las consultas masivas de los
  analistas tarden segundos y no horas.}
\lead{Automatización de cargas}{mantener los flujos de carga nocturna funcionando sin
  interrupciones ni errores humanos.}
\end{uxlist}

\personabloque{Pain points y desafíos}
\begin{uxlist}
\lead{Datos sin formato}{su mayor dolor de cabeza es cuando las áreas de negocio aprueban
  catálogos o envían información con formatos incorrectos, por ejemplo fechas en texto, que
  rompen los flujos de carga automatizados.}
\lead{Silos técnicos ocultos}{le molesta descubrir que un analista creó un almacén paralelo en
  su computadora, esquivando la arquitectura oficial del banco.}
\lead{Mantenimiento reactivo}{pierde mucho tiempo rastreando por qué un cálculo falló, solo
  para descubrir que alguien cambió una regla de negocio sin avisar a su equipo.}
\end{uxlist}

\personabloque{Comportamientos y hábitos}
\begin{uxlist}
\lead{Arquitectura primero}{antes de cargar un solo dato al sistema, dibuja diagramas de
  relación entre tablas.}
\lead{Validación estricta}{programa rutinas para que el sistema rechace automáticamente
  cualquier archivo que no cumpla con la arquitectura del almacén.}
\lead{Herramientas}{gestores de bases de datos masivas, herramientas de integración y software
  de modelado de datos institucionales.}
\end{uxlist}

\personabloque{Escenario de uso}
La Gerencia de Información aprueba una nueva métrica regulatoria. En lugar de levantar un
ticket complejo, Jorge entra al portal centralizado, revisa la definición exacta de la nueva
métrica en el catálogo de metadatos del negocio y mapea esa definición hacia una nueva tabla
en el almacén de datos. El portal actualiza el modelo sin interrumpir el servicio, permitiendo
que al día siguiente los analistas ya puedan consultar la métrica en sus tableros.

\subsection{Mapa de empatía de Jorge}
\mapaempatia
  {\item Ese archivo no pasó la validación de anoche por un error de formato.
   \item Por favor, no hagan consultas masivas directas a la tabla transaccional: usen las
     vistas del almacén.
   \item El diccionario de datos debe estar alineado a la estructura de la base.}
  {\item Los usuarios de negocio creen que subir un dato al sistema es magia; no ven la
     arquitectura que hay detrás.
   \item Necesitamos que todos usen el portal institucional para evitar los archivos fantasma.
   \item Cada cambio no anunciado es una falla que voy a rastrear de madrugada.}
  {\item Audita los tiempos de respuesta de la base de datos.
   \item Diseña arquitecturas para que la información histórica no ralentice el sistema actual.
   \item Rechaza integraciones que no cumplen con los estándares.}
  {\item Estrés al iniciar su jornada y ver que un proceso de carga nocturna falló.
   \item Frustración ante la falta de higiene de datos por parte de los originadores.
   \item Satisfacción técnica cuando una base de millones de registros responde en
     milisegundos.}

% ---------------------------------------------------------------- Persona 6
\clearpage
\section{Persona 6 · Arturo Castañeda}

\personahead{imagenes/persona_06_arturo.png}{Arturo Castañeda}
  {“En la Junta de Gobierno tengo cinco minutos para exponer; no me des bases de datos, dame
   certezas y tendencias claras.”}

\begin{personadatos}
Edad & 55 años \\
Sexo & Hombre \\
Ubicación & Ciudad de México \\
Ocupación & Director General de Estabilidad Financiera \\
Experiencia & 25 años en el sector público y financiero, toma de decisiones macroeconómicas \\
Nivel técnico & Básico: interpretación avanzada de gráficas y tableros, sin uso de SQL ni
  código \\
Modo principal & Directivo: abstracción, consolidación, visión estratégica e indicadores clave \\
\end{personadatos}

\personabloque{Antecedentes}
Arturo es un economista con amplia trayectoria. Como Director General, su responsabilidad es
monitorear los riesgos sistémicos del país y asesorar a los Subgobernadores y a la Gobernadora.
Su agenda está llena de reuniones de alto nivel. No tiene tiempo para limpiar datos ni cruzar
hojas de cálculo; consume la información ya procesada y consolidada. Está casado, tiene tres
hijos ya independientes y dos nietos. Fuera del trabajo lee ensayo económico, corre por las
mañanas y participa como profesor invitado en programas de posgrado.

\personabloque{Objetivos}
\begin{uxlist}
\lead{Visión macroscópica}{ver el estado general de la estabilidad financiera en un solo
  vistazo.}
\lead{Toma de decisiones rápidas}{identificar de inmediato anomalías, picos o tendencias
  preocupantes en el mercado.}
\lead{Certeza}{tener la absoluta seguridad de que el número que está viendo en la pantalla es
  la versión final, oficial y auditada por sus gerentes.}
\end{uxlist}

\personabloque{Pain points y desafíos}
\begin{uxlist}
\lead{Diferentes versiones de la verdad}{le incomoda mucho cuando en una reunión un director
  trae un reporte con un número y otro director trae un número distinto porque usaron bases de
  fechas diferentes.}
\lead{Falta de contexto inmediato}{a veces un indicador se ve alarmante y tiene que hacer
  varias llamadas a sus gerentes para entender por qué hubo ese pico.}
\lead{Ruido visual}{le mandan reportes con demasiadas tablas y datos crudos que no le ayudan a
  entender el panorama general.}
\end{uxlist}

\personabloque{Comportamientos y hábitos}
\begin{uxlist}
\lead{Consumo de información destilada}{su equipo le prepara resúmenes ejecutivos y tableros
  interactivos.}
\lead{Delegación profunda}{si encuentra una duda en un indicador, pide a sus gerentes que
  investiguen la raíz del problema en las bases subyacentes.}
\lead{Herramientas}{tableta institucional, correo electrónico, reportes ejecutivos en PDF y
  tableros interactivos.}
\end{uxlist}

\personabloque{Escenario de uso}
Quince minutos antes de entrar a presentar ante la Junta de Gobierno, Arturo abre el portal
desde su tableta y entra a su espacio directivo. Ve un tablero con los cuatro indicadores
clave de riesgo bancario del mes. Nota una ligera alza en uno de ellos. Toca la gráfica y el
portal le muestra una tarjeta de contexto escrita por el propietario del dato: “el alza se
debe al ajuste estacional aprobado la semana pasada”. Bloquea su tableta y entra a la reunión
con seguridad.

\subsection{Mapa de empatía de Arturo}
\mapaempatia
  {\item ¿Este dato ya está conciliado con la Dirección de Investigación Económica?
   \item Solo quiero ver la tendencia del trimestre, quiten el ruido de la gráfica.
   \item Si me preguntan en la Junta de dónde salió esta cifra, ¿qué respondo?}
  {\item Necesitamos una sola versión de la verdad en todo el Banco.
   \item No me importan los problemas técnicos; me importa tener la información para la
     reunión de las 12:00.
   \item El tiempo es mi recurso más escaso.}
  {\item Pide resúmenes ejecutivos en lugar de documentos de 50 páginas.
   \item Toma decisiones que pueden impactar a los mercados financieros basándose en la
     fiabilidad de estos reportes.
   \item Delega la validación técnica a sus gerentes y subgerentes.}
  {\item Mucha presión política e institucional por no dar un número equivocado.
   \item Frustración ante la burocracia que retrasa la disponibilidad de información.
   \item Seguridad plena cuando el sistema le demuestra que el dato es oficial.}

% ============================================================================
\clearpage
\aportacion{Arthur Jafed Zizumbo Velasco}{Matrícula A01796363 · Personas 7 y 8}

% ---------------------------------------------------------------- Persona 7
\section{Persona 7 · Mariana Ovalle Ríos}

\personahead{imagenes/persona_07_mariana.png}{Mariana Ovalle Ríos}
  {“Cada acceso que autorizo es una responsabilidad que firmo. Necesito darlo en minutos y
   poder retirarlo igual de rápido.”}

\begin{personadatos}
Edad & 41 años \\
Sexo & Mujer \\
Ubicación & Ciudad de México \\
Ocupación & Administradora de Plataformas de Información \\
Experiencia & 16 años en administración de sistemas institucionales y control de accesos \\
Nivel técnico & Avanzado: directorio corporativo, matrices de acceso, SQL y monitoreo de
  plataformas \\
Modo principal & Administración: altas y bajas de usuarios, asignación de roles y continuidad
  del servicio \\
\end{personadatos}

\personabloque{Antecedentes}
Mariana es ingeniera en sistemas computacionales con especialidad en seguridad de la
información. Coordina la administración de las plataformas de datos que utilizan varias
direcciones de la institución: da de alta a quien se incorpora, ajusta permisos cuando alguien
cambia de área y desactiva cuentas cuando alguien sale. Vive con su esposo y su hija
adolescente en el sur de la ciudad; fuera del trabajo corre medio maratones y toma cursos de
certificación en seguridad. Es metódica y desconfía de los procesos que dependen de que
alguien se acuerde.

\personabloque{Objetivos}
\begin{uxlist}
\lead{Altas y bajas oportunas}{dar acceso el primer día a quien se incorpora y retirarlo el
  mismo día de la baja, sin depender de recordatorios.}
\lead{Permiso mínimo suficiente}{asignar el rol más acotado que permita a cada persona hacer
  su trabajo, y poder demostrar por qué se otorgó.}
\lead{Continuidad del servicio}{anticipar saturación, fallas de sesión o interrupciones antes
  de que los usuarios las reporten.}
\end{uxlist}

\personabloque{Pain points y desafíos}
\begin{uxlist}
\lead{Accesos huérfanos}{cuentas que siguen activas después de que la persona cambió de área o
  dejó la institución, y que suelen detectarse hasta que alguien más las audita.}
\lead{Solicitudes sin contexto}{recibe peticiones de acceso por correo sin justificación, sin
  vigencia y sin visto bueno del propietario del dato.}
\lead{Trabajo invisible}{cuando la plataforma funciona nadie lo nota; cuando falla, todos
  escriben al mismo tiempo y no tiene con qué mostrar qué ocurrió.}
\end{uxlist}

\personabloque{Comportamientos y hábitos}
\begin{uxlist}
\lead{Verificación previa}{antes de otorgar un permiso confirma con el propietario del dato y
  deja registrada la autorización.}
\lead{Revisión periódica}{cruza la lista de usuarios activos contra la plantilla vigente para
  detectar cuentas que sobran.}
\lead{Herramientas}{consola de administración, directorio corporativo, hojas de control,
  sistema de tickets y tableros de monitoreo.}
\end{uxlist}

\personabloque{Escenario de uso}
Se incorpora una analista a la Dirección de Riesgos. Mariana entra al panel de administración
del portal, la da de alta con el rol Analista y el sistema le muestra exactamente qué fuentes
queda habilitada a consultar con ese rol. Al mes siguiente la persona se mueve a un área
directiva: Mariana cambia el rol en un clic, el portal reajusta de inmediato los módulos
visibles y deja registrado quién hizo el cambio y cuándo. Al cierre del trimestre descarga el
reporte de cuentas sin actividad y desactiva seis en una sola sesión, con la evidencia lista
por si auditoría la pide.

\subsection{Mapa de empatía de Mariana}
\mapaempatia
  {\item Necesito que el dueño del dato autorice por escrito, no por mensaje.
   \item ¿Esta cuenta sigue activa? La persona ya no está en esa área.
   \item Dame el rol correcto desde el inicio y nos ahorramos tres tickets.}
  {\item Si alguien entra donde no debe, la pregunta va a llegar a mi escritorio.
   \item Debería poder revisar todos los accesos en una pantalla, no en cinco sistemas.
   \item Cada excepción que hago hoy es un hallazgo de auditoría en seis meses.}
  {\item Da de alta, modifica y desactiva usuarios conforme cambia la plantilla.
   \item Cruza listas de acceso contra la plantilla para detectar cuentas huérfanas.
   \item Documenta cada autorización para poder sustentarla después.}
  {\item Responsabilidad constante por ser la última llave del acceso.
   \item Frustración cuando le piden accesos urgentes sin justificación.
   \item Tranquilidad cuando el sistema deja el rastro por ella.}

% ---------------------------------------------------------------- Persona 8
\clearpage
\section{Persona 8 · Ximena Solís Barrera}

\personahead{imagenes/persona_08_ximena.png}{Ximena Solís Barrera}
  {“No quiero descargar un archivo cada lunes. Dame un servicio estable, con su definición
   clara, y yo automatizo el resto.”}

\begin{personadatos}
Edad & 28 años \\
Sexo & Mujer \\
Ubicación & Ciudad de México, modalidad híbrida \\
Ocupación & Ingeniera de Aplicaciones Internas \\
Experiencia & 5 años desarrollando herramientas internas y servicios de datos \\
Nivel técnico & Avanzado: Python, APIs REST, control de versiones y automatización \\
Modo principal & Integración: consumo programático de datos y automatización de procesos \\
\end{personadatos}

\personabloque{Antecedentes}
Ximena es ingeniera en tecnologías de información y desarrolla las herramientas internas que
usan las áreas de negocio: calculadoras de riesgo, tableros a la medida y automatizaciones que
eliminan trabajo repetitivo. Comparte departamento con dos compañeras de la maestría, y dedica
los fines de semana a proyectos personales de programación y a andar en bicicleta. Le incomoda
construir sobre fuentes que pueden cambiar sin aviso, porque es ella quien recibe la llamada
cuando algo deja de funcionar.

\personabloque{Objetivos}
\begin{uxlist}
\lead{Contrato estable}{consumir datos mediante APIs con una definición documentada que no
  cambie sin aviso ni versión previa.}
\lead{Automatización confiable}{eliminar los pasos manuales de descarga y transformación que
  hoy repite cada semana.}
\lead{Autoservicio}{obtener credenciales, revisar el catálogo y probar una consulta sin abrir
  un ticket ni esperar días por una aprobación.}
\end{uxlist}

\personabloque{Pain points y desafíos}
\begin{uxlist}
\lead{Cambios silenciosos}{un campo cambia de nombre o de tipo y su proceso falla en producción
  sin que nadie lo hubiera anunciado.}
\lead{Documentación incompleta}{encuentra el servicio pero no la definición del campo, su
  periodicidad ni sus valores válidos.}
\lead{Fricción de acceso}{cada fuente nueva implica un trámite distinto, con formato,
  responsable y tiempo de respuesta diferentes.}
\end{uxlist}

\personabloque{Comportamientos y hábitos}
\begin{uxlist}
\lead{Prueba antes de integrar}{consulta el catálogo y ejecuta una muestra pequeña antes de
  escribir la automatización.}
\lead{Versiona todo}{guarda en el repositorio las consultas, los esquemas esperados y las
  pruebas que detectan cambios de estructura.}
\lead{Herramientas}{Python, clientes REST, control de versiones, contenedores y tableros de
  monitoreo.}
\end{uxlist}

\personabloque{Escenario de uso}
El área de riesgos le pide una herramienta que combine la exposición de créditos y derivados
por contraparte, actualizada cada mañana. Ximena entra al catálogo del portal, localiza los
dos conjuntos, revisa la definición y la periodicidad de cada campo y genera sus credenciales
desde su propio perfil. Prueba la consulta con un rango corto, confirma el formato de salida y
programa la ejecución diaria. Semanas después el portal anuncia un cambio de esquema y
mantiene disponible la versión anterior: su proceso sigue corriendo mientras ella migra sin
prisa.

\subsection{Mapa de empatía de Ximena}
\mapaempatia
  {\item ¿Dónde está la documentación de este servicio de datos?
   \item Si van a cambiar el esquema, avísenme antes, no cuando ya falló.
   \item Prefiero probar con una muestra antes de programar la carga completa.}
  {\item ¿Puedo confiar en que este campo signifique lo mismo el mes que viene?
   \item Automatizar esto ahorra cuatro horas a la semana; hacerlo mal cuesta el doble.
   \item Si tengo que pedir permiso cada vez, la gente lo va a resolver por fuera.}
  {\item Explora el catálogo y prueba consultas pequeñas antes de integrar.
   \item Versiona esquemas y escribe pruebas que detectan cambios de estructura.
   \item Automatiza descargas recurrentes y vigila que no fallen.}
  {\item Confianza cuando la definición del dato está documentada y versionada.
   \item Enojo cuando un cambio no anunciado rompe un proceso en producción.
   \item Satisfacción al eliminar un trabajo manual que alguien hacía cada semana.}

% ============================================================================
\clearpage
\section{Síntesis: persona primaria y alcance del método}

\subsection{Persona primaria}

Las ocho personas no se eligieron de una lista. Resultan de consolidar por objetivos de
trabajo los perfiles de la audiencia del apartado 1.1, agrupando en un mismo retrato a quienes
persiguen la misma meta aunque difieran en antecedentes; de ese conjunto se designa después
una persona primaria (Hartson y Pyla, 2012).

\begin{uxdestacado}
De las ocho personas, \textbf{\textcolor{uxnavy}{Laura Méndez}} (perfil operativo) se designa
como persona primaria del producto. Es el denominador común del portal: los siete perfiles
restantes también localizan y validan datos puntuales, aunque no sea su modo dominante.
Siguiendo el procedimiento de Hartson y Pyla (2012), la persona primaria es el blanco de
diseño específico y el criterio con el que se resuelven los conflictos de diseño: una interfaz
construida para Laura sigue sirviendo a las demás, mientras que una construida para Arturo
Castañeda, con su nivel de abstracción, no resolvería la tarea de Laura. Las siete personas
restantes se atienden sin sacrificar el caso primario.
\end{uxdestacado}

\subsection{Alcance y límite del método}

La literatura advierte que la técnica de personas rinde menos en dominios de trabajo
complejos, donde la práctica está muy normada y los individuos de un mismo rol resultan
intercambiables (Hartson y Pyla, 2012). El portal opera en uno de esos dominios. El equipo lo
asume y lo mitiga combinando las personas con una descripción explícita de perfiles y
responsabilidades sobre el dato (apartado 1.1), de modo que las decisiones de diseño no
dependan únicamente del retrato individual. Conviene además distinguir dos planos: los ocho
perfiles de la sección 1 son categorías de audiencia para la investigación, mientras que el
control de acceso del producto se resuelve con cuatro roles operativos: operativo, analista,
directivo y administrador.

\subsection{Criterios de calidad aplicados}

Las ocho personas se contrastaron con los criterios de calidad que propone la literatura:
deben ser ricas en detalle, relevantes para las decisiones de diseño, creíbles, específicas y
precisas, y deben evitar al usuario promedio, cuya mezcla de rasgos produce un retrato que no
representa a nadie (Hartson y Pyla, 2012). Por eso cada persona tiene nombre completo,
historia laboral, situación personal, herramientas concretas y un escenario de uso, en lugar
de una descripción genérica del puesto.

Existe además un riesgo particular en este proyecto: los integrantes del equipo trabajan en el
mismo dominio que se está diseñando, lo que facilita el sesgo de diseñar para uno mismo. La
literatura señala que las personas sirven precisamente para contrarrestar esa tendencia,
porque obligan a preguntarse cómo resolvería la tarea alguien con otras características y no
cómo la resolvería quien diseña (Hartson y Pyla, 2012). Las ocho personas cumplen aquí esa
función de control: obligan a evaluar cada decisión desde perfiles cuyas prioridades no
coinciden con las del equipo.

% ============================================================================
\clearpage
\section{Trazabilidad del problema al producto}

La siguiente matriz conecta las tres secciones de equipo con la sección individual: cada pain
point identificado en 1.3 corresponde a una necesidad de 1.2, a una característica del
producto de 3.2 y a la persona que la encarna. La cobertura es deliberada y completa: los ocho
perfiles de la audiencia tienen persona y característica asociada, y ninguna característica
del producto queda sin un pain point que la justifique.

\begin{uxtabla}{Z{0.85}L{2.5cm}Z{1.15}L{2.7cm}}{Trazabilidad de pain point a necesidad, característica del producto y persona}
\uxheadrow \thd{Pain point (1.3)} & \thd{Necesidad (1.2)} & \thd{Característica (3.2)} & \thd{Persona} \\ \midrule
Información dispersa & Localización & Buscador unificado y catálogo de datos & Laura Méndez \\
Fricción analítica & Profundidad & Explorador con filtros y exportación de datos crudos & Diego Hernández \\
Riesgo de interpretación & Confianza & Catálogo con definición, propietario y vigencia & Roberto Valdez \\
Evidencia limitada & Trazabilidad y control & Bitácora y evidencia & Elena Ruiz \\
Cambios e incidentes poco visibles & Operabilidad técnica & Monitoreo técnico & Jorge Mendieta \\
Silos de conocimiento & Confianza & Tableros consolidados con trazabilidad de cada cifra & Arturo Castañeda \\
Responsabilidad difusa & Gobierno & Acceso por permisos y roles con flujos de aprobación & Mariana Ovalle \\
Accesibilidad ineficiente & Interoperabilidad & Consumo por API con contrato documentado & Ximena Solís \\
\end{uxtabla}

% ============================================================================
\clearpage
\section{Referencias}

\begingroup
\setlength{\parskip}{9pt}
\setlength{\parindent}{-0.5cm}
\leftskip=0.5cm

de Voil, N. (2020). \textit{User experience foundations}. BCS, The Chartered Institute for IT.

Hartson, R., y Pyla, P. S. (2012). \textit{The UX book: Process and guidelines for ensuring a
quality user experience}. Morgan Kaufmann.

International Organization for Standardization. (2018). \textit{Ergonomics of human-system
interaction — Part 11: Usability: Definitions and concepts} (ISO 9241-11:2018).

Maioli, L. (2018). \textit{Fixing bad UX designs: Master proven approaches, tools, and
techniques to make your user experience great again}. Packt Publishing.

Osterwalder, A., y Pigneur, Y. (2010). \textit{Business model generation: A handbook for
visionaries, game changers, and challengers}. John Wiley \& Sons.

Portigal, S. (2013). \textit{Interviewing users}. Rosenfeld Media.

Ramírez Mejía, A. I. (2023). \textit{Grocery shopping app: A guided UX case study. Part 1.
Product and user definition} [Diapositivas]. TC4032 Experiencia del usuario y diseño de
interfaces, Maestría en Inteligencia Artificial Aplicada, Instituto Tecnológico y de Estudios
Superiores de Monterrey.

\endgroup
